\section*{Conventions and scope}
All spaces and coherent sheaves are complex analytic unless a polynomial ring is displayed explicitly. Cohomology of coherent sheaves has coefficients in $\C$. Integral and rational topological cohomology carry their indicated coefficients. Every compact complex manifold has its complex orientation.

We retain $f:X\to B=\PP^1_t$, $p_1=0$, $p_2=1$, $p_0=\infty$, and $m_1=3,m_2=4$. The smooth bielliptic reductions are $S_j=(f^{-1}(p_j))_{\rm red}$, with $f^*p_j=m_jS_j$. The fibre $W=f^{-1}(p_0)$ is reduced, irreducible and normal crossing, with normalization $\eta:\mathrm dP_6\to W$. The lattice $\Lambda=H_1(F,\Z)=\Z^4$ has the ordered basis $(\widehat\gamma,\widehat u,\widehat w,\widehat\delta)$. The monodromies $A_1,A_2$ have orders $3,4$; $M_0$ is unipotent with $(M_0-I)^2=0$. For the affine twists $v_1,v_2$ and integers $l_0,l_1,l_2$, we set $p=12l_0-4l_1-3l_2$ and retain $\pi_1(X)=\Z/|p|$.

On the upper half-plane with coordinate $z$, the generator fixed points are $z_1,z_2$ for $g_1,g_2\in\Delta(3,4,\infty)$. The marked period and its parameter conventions are
\[
 \Pi=\begin{pmatrix}6\mu&\tau&1&0\\ \beta&\mu&0&1\end{pmatrix},\qquad
 D=\Im\beta-\frac{6(\Im\mu)^2}{\Im\tau},\qquad
 \beta_c=\beta^{\rm part}+c.
\]
We write $M=\max_{B\setminus\{p_0\}}D_{\rm part}$; the parameter domain is $\Im c<-M$, and $w=e^{-\pi ic}$. Engel's rational elliptic surface $S\to\PP^1$ has fibres $\mathrm{IV}^*,\mathrm{III},\mathrm I_1$ at $0,1,\infty$, with Mordell--Weil generator $P$ and zero section $O$. Its line bundle is $M_S=\OO_S(P-O)$, its linearization scale is $u$, and its logarithmic-transformation data are $a_0,a_1$. The scalar $M$ and the line bundle $M_S$ retain these separate definitions throughout.

We distinguish integral cohomology $H^j$ from integral homology $H_j$; torsion in $H_2$ can coexist with $H^2=0$. The exponential is $z\mapsto e^{2\pi iz}$. A fibre without the word ``reduced'' means its full scheme-theoretic Cartier fibre. A normal-crossing support may have self-intersections and disconnected normalization. Engel's base labels $(0,1,\infty)$ correspond to $(p_1,p_2,p_0)$.

% Fragment for the first-page conventions. Exact checks: code/D_exact_checks.py.
Base meridians are counterclockwise in a local holomorphic coordinate,
including the coordinate $1/t$ at infinity. The analytic root deck
generators are clockwise: $s_j(g_jz)=e^{-2\pi i/m_j}s_j(z)$.
A positive base loop lifts to the inverse deck translate; identifying
its endpoint with the initial fibre applies $g_j$. The resulting
homology monodromies are the displayed $A_j$, with no inversion.
In Engel's common basis $\mathcal B_0=(e_1,e_2,e_3,e_4)$ the explicit
integral marking is
\[
P_{\rm E}=\begin{pmatrix}
0&-2&-1&-2\\2&-1&0&-1\\0&2&1&3\\-1&0&0&0
\end{pmatrix},\qquad \det P_{\rm E}=-1,
\qquad T_0P_{\rm E}=P_{\rm E}A_1,
\quad T_1P_{\rm E}=P_{\rm E}A_2.
\]
Its columns are the images of the retained ordered basis
$(\widehat\gamma,\widehat u,\widehat w,\widehat\delta)$.
Writing $\gamma=(1,0,0,0)$ and $\psi=(0,0,0,1)$, one has
$P_{\rm E}\widehat\delta=(-2,-1,3,0)^t=\delta_0$
and $\psi P_{\rm E}=-\gamma$.
Engel's positive-root affine generators are inverses of the retained
clockwise affine generators. Their invariant translations therefore
satisfy
\[
a_0=-P_{\rm E}v_1/3,\qquad a_1=-P_{\rm E}v_2/4.
\]
For $v_1=(1,2,-4,0)^t$ and $v_2=(-1,-3,3,0)^t$ this gives
$a_0=(0,0,0,1/3)^t$ and
$a_1=(-3,-1,3,-1)^t/4$.
Thus $12\psi(a_0+a_1)=1$, whereas the retained clockwise integer is
$p=-1$. With the transported inverse-period reference linearization,
$u=e^{-2\pi ic}=w^2$, $w=e^{-\pi ic}$, and $c\mapsto c+1$ changes
$(a_0,a_1)$ to $(a_0-\delta_0/3,a_1+\delta_0/4)$.

